\chapter{Analytical design}
\label{ch:analytical}

This chapter fixes the operating point and the passive components from
closed-form relations. Everything here is computed before any simulation, so
that the simulation of \cref{ch:model} has something independent to be checked
against.

\section{Operating point}

With space-vector modulation the maximum phase voltage amplitude that can be
synthesised without overmodulation is
%
\begin{equation}
  \hat{V}_{\mathrm{ph,max}} = \frac{\Vdc}{\sqrt{3}}
  = \frac{\SI{800}{\volt}}{\sqrt{3}} = \SI{462}{\volt} .
  \label{eq:svpwm-limit}
\end{equation}
%
A phase voltage amplitude of $\hat{V}_\mathrm{ph} = \SI{427}{\volt}$ is
chosen, giving a modulation index
%
\begin{equation}
  m = \frac{\hat{V}_\mathrm{ph}}{\Vdc} = \frac{427}{800} = 0.534 ,
  \qquad m_{\max} = \frac{1}{\sqrt{3}} = 0.577 ,
\end{equation}
%
that is, \SI{92}{\percent} of the linear range. The remaining headroom is
deliberate: it leaves room for the voltage lost to the dead time and to the
device drops, and it keeps the converter out of overmodulation during
transients, where the harmonic content degrades sharply.

The resulting load is an equivalent $RL$ branch per phase drawing
\SI{30}{\Arms}, with $R = \SI{9.43}{\ohm}$ and
$L = \SI{2.5}{\milli\henry}$.

\subsection{Power factor and which power is quoted}
\label{sec:powerfactor}

That branch is \emph{not} at unity power factor, and the distinction matters
for every figure derived from it. At the \SI{200}{\hertz} fundamental
%
\begin{equation}
  X_L = 2\pi f_\mathrm{ref} L = \SI{3.142}{\ohm},
  \qquad
  |Z| = \sqrt{R^2 + X_L^2} = \SI{9.940}{\ohm},
  \qquad
  \cos\varphi = \frac{R}{|Z|} = 0.949 ,
  \label{eq:powerfactor}
\end{equation}
%
that is, a displacement angle of \SI{18.4}{\degree}. The values of $R$ and
$L$ are the ones that give \SI{30}{\Arms} at the chosen modulation depth;
it is only the unity-power-factor label that would be wrong.

Three different powers can be quoted for this operating point, and they
differ by more than a rounding error:
%
\begin{align}
  S_1 &= 3\,V_{1,\mathrm{rms}} I_{1,\mathrm{rms}}
       = 3 \cdot \frac{427}{\sqrt{2}} \cdot 30
       = \SI{27.17}{\kilo\VA}
       &&\text{(apparent, fundamental)} \label{eq:apparent}\\
  P   &= S_1 \cos\varphi = \SI{25.79}{\kilo\watt}
       &&\text{(active, ideal)} \label{eq:active-ideal}\\
  P   &= R\left(I_1^2 + I_2^2 + I_3^2\right) = \SI{25.48}{\kilo\watt}
       &&\text{(active, simulated)} \label{eq:active-sim}
\end{align}

This report quotes the active power, \SI{25.48}{\kilo\watt}, because
that is the quantity against which an efficiency is meaningful. The value
comes from \cref{eq:active-sim} rather than from
\cref{eq:active-ideal}: on a series $RL$ branch the inductance dissipates
nothing, so all the active power appears in $R$, and the relation is exact
without needing a voltage measurement. The \SI{1.2}{\percent} gap between
the two is the voltage lost to the \SI{260}{\nano\second} dead time
(\cref{sec:deadtime}), which reduces the fundamental from \SI{301.9}{} to
\SI{298.3}{\Vrms}; the modulation amplitude of \SI{427}{\volt} is the one
that brings the current back to \SI{30}{\Arms} with that dead time.

\section{Current ratings}

The peak phase current in the ideal sinusoidal case is
$\hat{I} = \sqrt{2}\cdot 30 = \SI{42.4}{\ampere}$. The simulated waveform, which
carries the switching ripple on top of the fundamental, peaks at
\SI{43.0}{\Apk}; this is the value the devices see and the one that sizes the overload
protection of the gate drive.

The RMS current of each device follows from synchronous rectification: the
channel conducts in both directions whenever the gate is on, so the upper
switch of a leg carries $i^2$ for a fraction $d(t) = \tfrac12[1 + m\sin\omega t]$
of each period. Averaged over the fundamental, the term in $m$ multiplies an
odd function and vanishes, leaving
%
\begin{equation}
  I_\mathrm{rms,dev} = \sqrt{\overline{d(t)\,i^2(t)}}
  = \frac{\hat{I}}{2} = \frac{I_\mathrm{ph}}{\sqrt{2}} = \SI{21.2}{\ampere},
\end{equation}
%
independent of modulation index and power factor, dead time neglected. It is
the current that sets the conduction loss, $P_\mathrm{cond} = \Rdson\,
I_\mathrm{rms,dev}^2$ per device.

Against the device rating this is comfortable: the \device{} is rated well
above \SI{42}{\ampere} continuous at the relevant case temperature. The
limiting quantity is the junction temperature, which \cref{ch:thermal}
resolves.

\section{DC-link capacitor}
\label{sec:dclink}

\subsection{Sizing for ripple}

The capacitor must hold the bus within \SI{0.5}{\percent} of
\SI{800}{\volt}, i.e. $\Delta V \leq \SI{4}{\volt}$. The worst-case charge
displaced in one switching period is bounded by
%
\begin{equation}
  \Delta Q \approx \frac{\hat{I}}{2}\cdot\frac{1}{2\fsw},
  \qquad
  C \geq \frac{\Delta Q}{\Delta V},
\end{equation}
%
which for $\hat{I}=\SI{42.4}{\ampere}$, $\fsw=\SI{20}{\kilo\hertz}$ and
$\Delta V = \SI{4}{\volt}$ gives $C \gtrsim \SI{133}{\micro\farad}$ as an
order of magnitude.

\begin{finding}{An ideal source hid the ripple that the hand calculation predicted}
The first version adopted \SI{100}{\micro\farad}, reassured by a simulation
that gave \SI{0.43}{\volt} of ripple, an order of magnitude inside the
requirement. That model fed the DC-link from an ideal voltage source, which
absorbed half of the ripple current in place of the capacitor. With a
realistic source inductance of \SI{5}{\micro\henry} the ripple with
\SI{100}{\micro\farad} reaches \SI{6.2}{\volt} in the worst case, 54 percent
beyond the limit: the model's defect was hiding exactly the violation that
the \SI{133}{\micro\farad} estimate above predicted. With the corrected
source, hand calculation and simulation agree.
\end{finding}

$C_\mathrm{DC} = \SI{200}{\micro\farad}$ is adopted, as two film capacitors
of \SI{140}{} and \SI{60}{\micro\farad} in the \SI{1100}{\volt}
class~\cite{vishay_mkp1848}. The ripple is \SI{1.36}{\volt} at the rated
point (\SI{0.17}{\percent}) and \SI{2.86}{\volt} in the worst case, inside
the \SI{4}{\volt} limit. Two further conditions come with it. A source that
is both less inductive than \SI{1.5}{\micro\henry} and less damped than
\SI{50}{\milli\ohm} at \SI{20}{\kilo\hertz} would tune the
source--capacitor resonance onto the switching harmonics, and no reasonable
capacitor closes that band: it is an assumption on the battery and harness
(\cref{tab:sysassume}). And the film's permitted voltage falls with its
hotspot temperature (\SI{1100}{\volt} at \SI{85}{\celsius},
\SI{800}{\volt} at \SI{105}{\celsius}), so the \SI{1100}{\volt} class, not
the \SI{900}{\volt} one, is what allows \SI{900}{\volt} continuous up to a
hotspot of \SI{98}{\celsius}.

\subsection{Sizing for RMS current}

\begin{figure}[htbp]
\centering
% Chi porta cosa nel DC-link: la corrente del ponte e' a gradini, quella della batteria e'
% continua, il condensatore porta la differenza. Forma qualitativa su tre periodi di commutazione.
\begin{tikzpicture}
  \pgfplotsset{dcax/.style={width=0.82\textwidth, height=3.6cm, xmin=0, xmax=3,
    xtick={0,1,2,3}, grid=major, grid style={black!10},
    tick label style={font=\scriptsize}, label style={font=\footnotesize},
    legend style={font=\scriptsize, draw=none, fill=white, fill opacity=0.85, text opacity=1,
                  at={(1.0,1.0)}, anchor=north east, legend columns=2}}}
  \begin{axis}[dcax, name=sopra, ymin=-5, ymax=45, ytick={0,20,40}, xticklabels={},
               ylabel={current [A]}]
    % zero vettori: nessuna corrente; vettori attivi: la corrente di una fase o di due
    \addplot[farancio, thick, const plot] coordinates {
      (0,0) (0.15,36) (0.35,10) (0.65,36) (0.85,0)
      (1,0) (1.15,36) (1.35,10) (1.65,36) (1.85,0)
      (2,0) (2.15,36) (2.35,10) (2.65,36) (2.85,0) (3,0)};
    \addlegendentry{bridge input $i_\mathrm{in}$}
    \addplot[fblu, very thick, domain=0:3] {17.4};
    \addlegendentry{battery $I_\mathrm{dc}$, its average}
  \end{axis}
  \begin{axis}[dcax, at={(sopra.south)}, anchor=north, yshift=-3mm, ymin=-25, ymax=25,
               ytick={-20,0,20}, xlabel={time [switching periods]}, ylabel={current [A]}]
    \addplot[fgrigio, thick, const plot] coordinates {
      (0,-17.4) (0.15,18.6) (0.35,-7.4) (0.65,18.6) (0.85,-17.4)
      (1,-17.4) (1.15,18.6) (1.35,-7.4) (1.65,18.6) (1.85,-17.4)
      (2,-17.4) (2.15,18.6) (2.35,-7.4) (2.65,18.6) (2.85,-17.4) (3,-17.4)};
    \addlegendentry{capacitor $i_C = i_\mathrm{in} - I_\mathrm{dc}$, zero average}
  \end{axis}
\end{tikzpicture}

\caption{Who carries what in the DC-link, over three switching periods (shape only).
The bridge draws its current in blocks, one phase current or the sum of two
depending on the active vector; the battery, behind its cable inductance,
supplies only the average; the capacitor supplies the rest, with zero
average and an RMS value that sets its heating.}
\label{fig:dccurrent}
\end{figure}

The capacitor carries the difference between the inverter input current and
the (smooth) DC source current (\cref{fig:dccurrent}). Its RMS value is the quantity that determines
the capacitor's self-heating and therefore its life. A closed-form result for
the two-level bridge is given by Kolar and Round~\cite{kolar_dclink}:
%
\begin{equation}
  \frac{I_{C,\mathrm{rms}}}{I_\mathrm{out,rms}}
  = \sqrt{ 2m\left[\frac{\sqrt{3}}{4\pi}
    + \cos^2\!\varphi\left(\frac{\sqrt{3}}{\pi}-\frac{9}{16}m\right)\right]} .
  \label{eq:kolar}
\end{equation}

\begin{finding}{Two different modulation indices live in this report}
\Cref{eq:kolar} uses $m$ referred to \emph{half} the DC-link,
$m = 2\hat{V}_\mathrm{ph}/\Vdc = 1.07$, which is the convention of the
original derivation. \Cref{eq:svpwm-limit} uses $m$ referred to the full
DC-link, $m = \hat{V}_\mathrm{ph}/\Vdc = 0.534$.

The distinction matters: with $m = 1.07$ the bracket
$\left(\sqrt{3}/\pi - 9m/16\right)$ is \emph{negative}, so at the rated point
the ratio grows as $\cos\varphi$ falls; with $m = 0.534$ it is positive and
the trend reverses. Substituting the wrong one places the rated point on the
wrong branch of the curve.
\end{finding}

Two properties of \cref{eq:kolar} matter for the design. At the \emph{rated}
modulation index, $m = 1.07$, the bracket multiplying $\cos^2\varphi$ is
negative and the ratio grows as $\cos\varphi$ falls: \num{0.447} at the rated
point, \num{0.542} at $\cos\varphi \to 0$. But the modulation index
falls with speed, and the worst case must be sought over the whole
operating envelope, $0 \leq m \leq 1.15$ and $0 \leq \cos\varphi \leq 1$.
There the maximum is at unity power factor, where
$\partial/\partial m$ of the bracket vanishes:
%
\begin{equation}
  m^\ast = \frac{8}{9}\left(\frac{\sqrt{3}}{4\pi}+\frac{\sqrt{3}}{\pi}\right)
         = 0.613,
  \qquad
  \frac{I_{C,\mathrm{rms}}}{I_\mathrm{out,rms}}\bigg|_{\max} = 0.650 ,
\end{equation}
%
that is \SI{19.5}{\ampere} at \SI{30}{\Arms} (\cref{fig:kolar}). The capacitor's worst case is a
low-speed point at full current and unity power factor, not the rated point
(\SI{13.4}{\ampere}) nor the reactive one (\SI{16.9}{\ampere} at $m = 1.15$).

\begin{figure}[htbp]
\centering
% Corrente efficace nel DC-link riferita alla corrente di fase, al variare di
% M_f = V_M/V_DC e del fattore di potenza. Legenda sotto il grafico, etichette
% dei due punti fuori dalle curve.
\begin{tikzpicture}
  \begin{axis}[
    width=0.82\textwidth, height=6.2cm,
    xmin=0, xmax=0.6, ymin=0, ymax=0.8,
    xtick={0,0.1,0.2,0.3,0.4,0.5,0.6},
    xticklabel style={/pgf/number format/fixed, /pgf/number format/precision=1},
    scaled x ticks=false,
    xlabel={modulation index $M_f = V_M/\Vdc$},
    ylabel={$I_{C,\mathrm{RMS}}/I_{o,\mathrm{RMS}}$},
    grid=major, grid style={black!10},
    legend columns=4,
    legend style={font=\scriptsize, draw=none, fill=none, at={(0.5,-0.32)}, anchor=north,
                  /tikz/every even column/.append style={column sep=6pt}},
    tick label style={font=\scriptsize}, label style={font=\footnotesize},
    domain=0.001:0.577, samples=120]
    \addplot[fblu, thick] {sqrt(max(0, x*(sqrt(3)/pi + 1*(4*sqrt(3)/pi - 4.5*x))))};
    \addlegendentry{$\cos\varphi = 1$}
    \addplot[farancio, thick] {sqrt(max(0, x*(sqrt(3)/pi + 0.949^2*(4*sqrt(3)/pi - 4.5*x))))};
    \addlegendentry{$\cos\varphi = 0.949$ (rated)}
    \addplot[fverde, thick, dashed] {sqrt(max(0, x*(sqrt(3)/pi + 0.25*(4*sqrt(3)/pi - 4.5*x))))};
    \addlegendentry{$\cos\varphi = 0.5$}
    \addplot[fgrigio, thick, dotted] {sqrt(max(0, x*(sqrt(3)/pi)))};
    \addlegendentry{$\cos\varphi = 0$}
    % limite lineare SVPWM
    \draw[black!50, dashed] (axis cs:0.577,0) -- (axis cs:0.577,0.8);
    \node[font=\scriptsize, black!60, anchor=south east] at (axis cs:0.577,0.02) {$1/\sqrt{3}$};
    % punto nominale e caso peggiore, con le etichette su fondo bianco
    \addplot[only marks, mark=*, mark size=1.8pt, farancio] coordinates {(0.534,0.447)};
    \node[font=\scriptsize, farancio, fill=white, inner sep=1pt, anchor=north] at (axis cs:0.47,0.36) {rated point, 0.447};
    \draw[farancio, thin] (axis cs:0.534,0.447) -- (axis cs:0.49,0.365);
    \addplot[only marks, mark=*, mark size=1.8pt, fblu] coordinates {(0.306,0.650)};
    \node[font=\scriptsize, fblu, fill=white, inner sep=1pt, anchor=south] at (axis cs:0.306,0.70) {worst case, 0.650 at $M_f = 0.306$};
    \draw[fblu, thin] (axis cs:0.306,0.650) -- (axis cs:0.306,0.70);
  \end{axis}
\end{tikzpicture}

\caption{DC-link RMS current relative to the phase current, \cref{eq:kolar}, over
the whole modulation range and four power factors. The rated point sits on
the falling branch of its curve; the worst case of the envelope is a
low-speed point at unity power factor.}
\label{fig:kolar}
\end{figure}

A rule of thumb frequently used in practice sets
$I_{C,\mathrm{rms}} \approx 0.65\,I_\mathrm{out,rms}$. It coincides with the
maximum of \cref{eq:kolar} over the
envelope, which is presumably where it comes from. \Cref{ch:model} checks the
formula against the simulation at the rated modulation index.

The two film capacitors are rated \SI{43.5}{} and \SI{40}{\ampere} RMS at
\SI{85}{\celsius}~\cite{vishay_mkp1848}. At the worst case the larger one
carries about \SI{13.6}{\ampere} of the \SI{19.5}{}, which heats its hotspot
by \SI{1.5}{\celsius}: the binding constraint is the voltage class. Film is chosen over electrolytic for the ripple-current capability
and the life at temperature; the cost is volume, \SI{864}{\gram} and two
packages of \SI{130}{} and \SI{65}{\milli\metre} by \SI{65}{\milli\metre}.

\section{Switching frequency}

The choice of \SI{20}{\kilo\hertz} balances three effects that pull in
different directions. At \SI{20}{\kilo\hertz} the acoustic noise from
magnetostriction in the machine and from the capacitor is above the range
most people hear; \SI{10}{\kilo\hertz} would halve the switching losses but
put a tone in the middle of the audible band, which a passenger vehicle does
not accept. A higher $\fsw$ reduces the current ripple and with it the
harmonic losses in the machine. The switching losses grow linearly with
$\fsw$: \cref{ch:losses} measures the slope at \SI{596}{\micro\joule} per
device per cycle.

\Cref{ch:losses} shows the full sweep from \SI{10}{} to \SI{30}{\kilo\hertz}.
The total loss varies from \SI{90.5}{\watt} to \SI{166.1}{\watt} across
that range, and the junction temperature at the rated point from
\SI{76.9}{\celsius} to \SI{86.8}{\celsius}: at the rated point the design is
not thermally limited anywhere in the range. The limit is at locked rotor
(\cref{sec:lockedrotor}), where the \SI{20}{\kilo\hertz} is held by
requirement and costs its full switching loss on one device. The frequency
is therefore chosen on the acoustic and harmonic-content arguments and then
verified at locked rotor.
